\section{实验}
\label{sec:experiments-zh}

\subsection{实验设置}

\subsubsection{数据集}

本文在 ACM、DBLP、AMiner 和 IMDB 四个异构图基准数据集上评估 DVCL。每个数据集包含多种节点类型和关系类型，任务是在目标节点类型上进行半监督分类。按照 HSeCo 的协议，本文使用数据集特定的元路径构造语义邻域。表~\ref{tab:dataset_settings_zh} 总结了目标节点类型和元路径。最终主对比将至少包含三个具有完整五随机种子结果的数据集；第四个数据集将在完成相同协议后报告。

\begin{table}[htbp]
\centering
\caption{数据集设置。}
\label{tab:dataset_settings_zh}
\begin{tabular}{lll}
\toprule
数据集 & 目标节点类型 & 元路径 \\
\midrule
ACM & Paper & P-A-P, P-F-P \\
DBLP & Author & A-P-A, A-P-C-P-A, A-P-T-P-A \\
AMiner & Paper & P-A-P, P-R-P \\
IMDB & Movie & M-D-M, M-A-M \\
\bottomrule
\end{tabular}
\end{table}

\subsubsection{评估协议}

本文只使用分类准确率作为评估指标：
\begin{equation}
\operatorname{Accuracy}
=\frac{1}{|\mathcal{V}_{\mathrm{eval}}|}
\sum_{i\in\mathcal{V}_{\mathrm{eval}}}
\mathbb{I}(\hat{y}_i=y_i).
\end{equation}
除非特别说明，每个表格报告随机种子 1、2、3、4、5 的均值和标准差。对于每个随机种子，所有方法使用完全相同的训练集、验证集和测试集划分，以及相同的攻击图、目标节点 ID 和评估节点。超参数只能依据验证集选择，并为所有方法分配相同的调参预算。表格中的“--”表示该结果尚未在统一协议下产生。当存在配对的 HSeCo 参照结果时，准确率下方的小号彩色文字表示相对同列 HSeCo 行的百分点变化；绿色和“+”表示上升，红色和“-”表示下降。预备表使用明确标注的审计前 HSeCo 行，最终表则使用 HSeCo$^{\dagger}$。

基线原论文中的数字不会与本文重新运行的结果混放在主表中。HSeCo$^{\dagger}$ 表示本文审计后的复现版本。HSeCo 原论文结果与本文复现结果之间的差异将在附录中单独说明，并同时给出数据划分、配置和 checkpoint 行为。

\subsubsection{基线方法}

表~\ref{tab:baselines_zh} 给出统一比较所使用的基线，覆盖标准异构图架构、异构图对比学习方法和面向结构扰动的防御方法。GCN-Jaccard、GNNGuard 等同构图防御方法暂不进入主比较；只有在为所有此类方法预先固定同一种异构图到同构图的转换方式后，才能作为补充基线加入。

\begin{table}[htbp]
\centering
\small
\caption{统一实验中包含的方法。}
\label{tab:baselines_zh}
\begin{tabular}{lp{0.21\linewidth}p{0.56\linewidth}}
\toprule
方法 & 类别 & 在比较中的作用 \\
\midrule
HAN & 异构图基础模型 & 使用元路径注意力、但没有显式鲁棒机制的标准基线。 \\
HGT & 异构图基础模型 & 使用类型感知投影和注意力的 Transformer 基线。 \\
MAGNN & 异构图基础模型 & 使用细粒度元路径实例聚合的语义建模基线。 \\
HeCo & 对比学习 HGNN & 不包含显式攻击防御的异构图对比学习基线。 \\
HSeCo$^{\dagger}$ & 鲁棒 HGNN & 本文在统一协议下重新运行的直接语义拓扑基线。 \\
RoHe & 鲁棒 HGNN & 面向对抗性结构扰动的异构图防御基线。 \\
HeteGuard & 鲁棒 HGNN & 抑制不可靠结构信息的异构图防御基线。 \\
DVCL & 本文方法 & 结合净化语义拓扑、特征诱导图和跨视图对比对齐。 \\
\bottomrule
\end{tabular}
\end{table}

\subsubsection{攻击设置}

对于全局攻击，本文评估 PRBCD 和 HetePRBCD，扰动率为 $5\%$、$10\%$、$15\%$、$20\%$ 和 $25\%$。PRBCD 执行全局结构扰动，HetePRBCD 在生成攻击时保留异构关系约束。表中同时报告干净图准确率和受攻击准确率。

对于目标攻击，本文使用 FGSM 风格的异构图结构攻击。在 ACM 上，每个随机种子从测试集中选取 500 个目标 paper 节点，所有方法复用相同的目标节点 ID。最终目标攻击表报告 $\Delta\in\{1,2,3\}$ 三种预算，准确率只在被选中的目标节点上计算。

当前攻击主要修改图拓扑，不修改节点属性。由于 DVCL 的一个视图由节点特征构造，本文还预留了特征扰动以及同时针对两个视图的联合自适应攻击实验。在这些实验完成之前，论文的经验结论严格限定为对结构扰动的鲁棒性。

\subsubsection{实现与模型选择}

最大训练轮数为 200，early stopping 依据验证集性能执行，优化器使用 Adam。每个 checkpoint 必须包含推理所需的全部可训练模块。对于 HSeCo$^{\dagger}$，checkpoint 同时保存并恢复 HAN 语义模块和节点分类器；恢复 checkpoint 后，再使用已恢复的语义模块重新生成语义图。

DVCL 的特征诱导视图使用原始目标节点特征、有向 KNN 和 $k=20$。拓扑视图沿用 HSeCo 的语义拓扑构造，两个视图的表示通过拼接用于分类。本文设置 $\lambda_{\mathrm{HAN}}=1.0$。根据表~\ref{tab:param_zh} 中的预备参数实验，统一重跑前固定 $\lambda_{\mathrm{DVCL}}=1.0$。该设置兼顾干净准确率和受攻击准确率，之后不再针对不同数据集或测试攻击单独调参。

\FloatBarrier
\subsection{主要结果}

\subsubsection{跨数据集汇总}

表~\ref{tab:cross_dataset_zh} 是统一重跑后需要补齐的最终汇总表。“攻击平均”表示 PRBCD 和 HetePRBCD 共十种攻击条件的算术平均。当前已经填入的 ACM 数值属于预备记录，用于完善论文结构，但不作为最终外部基线比较。

\begin{table}[htbp]
\centering
\small
\setlength{\tabcolsep}{3.5pt}
\caption{统一协议下的跨数据集准确率。最终结果报告五个随机种子的均值 $\pm$ 标准差。}
\label{tab:cross_dataset_zh}
\resizebox{\textwidth}{!}{%
\begin{tabular}{lcccccccc}
\toprule
& \multicolumn{2}{c}{ACM} & \multicolumn{2}{c}{DBLP} & \multicolumn{2}{c}{AMiner} & \multicolumn{2}{c}{IMDB} \\
\cmidrule(lr){2-3}\cmidrule(lr){4-5}\cmidrule(lr){6-7}\cmidrule(lr){8-9}
方法 & 干净 & 攻击平均 & 干净 & 攻击平均 & 干净 & 攻击平均 & 干净 & 攻击平均 \\
\midrule
HAN & -- & -- & -- & -- & -- & -- & -- & -- \\
HGT & -- & -- & -- & -- & -- & -- & -- & -- \\
MAGNN & -- & -- & -- & -- & -- & -- & -- & -- \\
HeCo & -- & -- & -- & -- & -- & -- & -- & -- \\
HSeCo$^{\dagger}$ & -- & -- & -- & -- & -- & -- & -- & -- \\
RoHe & -- & -- & -- & -- & -- & -- & -- & -- \\
HeteGuard & -- & -- & -- & -- & -- & -- & -- & -- \\
DVCL ($\lambda=1.0$，预备) & $0.8810\pm0.0049$ & 0.8723 & -- & -- & -- & -- & -- & -- \\
DVCL ($\lambda=1.0$，最终) & -- & -- & -- & -- & -- & -- & -- & -- \\
\bottomrule
\end{tabular}}
\end{table}

\subsubsection{ACM 全局结构攻击}

表~\ref{tab:acm_prbcd_zh} 和表~\ref{tab:acm_heteprbcd_zh} 将模型放在行中，后续增加基线时无需修改表格结构。审计前 HSeCo 行和 DVCL（$\lambda=0.5$）行保留已有历史记录，彩色变化在同一列内配对计算。所有模型在相同冻结攻击产物上完成后，这两行将由 HSeCo$^{\dagger}$ 和固定 $\lambda=1.0$ 的 DVCL 配置替代。

\begin{table}[htbp]
\centering
\small
\caption{PRBCD 下 ACM 的准确率。当前数值为预备五随机种子记录。}
\label{tab:acm_prbcd_zh}
\resizebox{\textwidth}{!}{%
\begin{tabular}{lcccccc}
\toprule
方法 & 干净 & 5\% & 10\% & 15\% & 20\% & 25\% \\
\midrule
HAN & -- & -- & -- & -- & -- & -- \\
HGT & -- & -- & -- & -- & -- & -- \\
MAGNN & -- & -- & -- & -- & -- & -- \\
HeCo & -- & -- & -- & -- & -- & -- \\
HSeCo（审计前） & $0.8679\pm0.0129$ & $0.8491\pm0.0086$ & $0.8519\pm0.0080$ & $0.8484\pm0.0091$ & $0.8530\pm0.0092$ & $0.8504\pm0.0092$ \\
HSeCo$^{\dagger}$（最终） & -- & -- & -- & -- & -- & -- \\
RoHe & -- & -- & -- & -- & -- & -- \\
HeteGuard & -- & -- & -- & -- & -- & -- \\
DVCL ($\lambda=0.5$，预备) & \accgain{$0.8770\pm0.0071$}{+0.91 pp} & \accgain{$0.8732\pm0.0092$}{+2.41 pp} & \accgain{$0.8718\pm0.0077$}{+1.99 pp} & \accgain{$0.8674\pm0.0060$}{+1.90 pp} & \accgain{$0.8660\pm0.0146$}{+1.30 pp} & \accgain{$0.8659\pm0.0082$}{+1.55 pp} \\
DVCL ($\lambda=1.0$，最终) & -- & -- & -- & -- & -- & -- \\
\bottomrule
\end{tabular}}
\end{table}

\begin{table}[htbp]
\centering
\small
\caption{HetePRBCD 下 ACM 的准确率。当前数值为预备五随机种子记录。}
\label{tab:acm_heteprbcd_zh}
\resizebox{\textwidth}{!}{%
\begin{tabular}{lcccccc}
\toprule
方法 & 干净 & 5\% & 10\% & 15\% & 20\% & 25\% \\
\midrule
HAN & -- & -- & -- & -- & -- & -- \\
HGT & -- & -- & -- & -- & -- & -- \\
MAGNN & -- & -- & -- & -- & -- & -- \\
HeCo & -- & -- & -- & -- & -- & -- \\
HSeCo（审计前） & $0.8679\pm0.0129$ & $0.8530\pm0.0121$ & $0.8589\pm0.0126$ & $0.8546\pm0.0070$ & $0.8548\pm0.0124$ & $0.8511\pm0.0092$ \\
HSeCo$^{\dagger}$（最终） & -- & -- & -- & -- & -- & -- \\
RoHe & -- & -- & -- & -- & -- & -- \\
HeteGuard & -- & -- & -- & -- & -- & -- \\
DVCL ($\lambda=0.5$，预备) & \accgain{$0.8770\pm0.0071$}{+0.91 pp} & \accgain{$0.8762\pm0.0027$}{+2.33 pp} & \accgain{$0.8737\pm0.0087$}{+1.48 pp} & \accgain{$0.8726\pm0.0046$}{+1.79 pp} & \accgain{$0.8700\pm0.0108$}{+1.53 pp} & \accgain{$0.8703\pm0.0063$}{+1.92 pp} \\
DVCL ($\lambda=1.0$，最终) & -- & -- & -- & -- & -- & -- \\
\bottomrule
\end{tabular}}
\end{table}

现有 ACM 记录表明，预备 DVCL 配置在已测试的结构扰动率下比审计前的 HSeCo 复现更稳定。这一观察仍然是临时结论；最终结论只依据冻结实验协议以及固定 $\lambda_{\mathrm{DVCL}}=1.0$ 后重新产生的结果。

\subsubsection{目标结构攻击}

表~\ref{tab:target_attack_zh} 按三种攻击预算报告目标节点准确率。目前只有 HSeCo 和预备 DVCL 配置在 $\Delta=3$ 下具有记录。预算 1、2，外部基线以及最终 DVCL 配置仍需在相同目标节点 ID 上补齐。

\begin{table}[htbp]
\centering
\small
\caption{FGSM 风格结构攻击下 ACM 目标节点准确率。}
\label{tab:target_attack_zh}
\begin{tabular}{lcccc}
\toprule
方法 & 干净 & $\Delta=1$ & $\Delta=2$ & $\Delta=3$ \\
\midrule
HAN & -- & -- & -- & -- \\
HGT & -- & -- & -- & -- \\
MAGNN & -- & -- & -- & -- \\
HeCo & -- & -- & -- & -- \\
HSeCo（审计前） & 0.8624 & -- & -- & 0.7608 \\
HSeCo$^{\dagger}$（最终） & -- & -- & -- & -- \\
RoHe & -- & -- & -- & -- \\
HeteGuard & -- & -- & -- & -- \\
DVCL ($\lambda=0.5$，预备) & \accgain{0.8724}{+1.00 pp} & -- & -- & \accgain{0.8608}{+10.00 pp} \\
DVCL ($\lambda=1.0$，最终) & -- & -- & -- & -- \\
\bottomrule
\end{tabular}
\end{table}

\FloatBarrier
\subsection{消融与敏感性分析}

\subsubsection{对比损失权重}

表~\ref{tab:param_zh} 总结了已有的 $\lambda_{\mathrm{DVCL}}$ 参数实验，其中攻击平均覆盖十种全局攻击条件。本文选择 $\lambda_{\mathrm{DVCL}}=1.0$，因为它取得最高的干净准确率，同时受攻击准确率接近最优，适合作为统一重跑的平衡配置。本文不会仅因为 2.0 的攻击平均略高而选择该值。

\begin{table}[htbp]
\centering
\caption{ACM 上 $\lambda_{\mathrm{DVCL}}$ 的准确率敏感性。}
\label{tab:param_zh}
\begin{tabular}{ccc}
\toprule
$\lambda_{\mathrm{DVCL}}$ & 干净准确率 & 攻击平均准确率 \\
\midrule
0.1 & -- & 0.8665 \\
0.5 & $0.8770\pm0.0071$ & 0.8710 \\
1.0 & $0.8810\pm0.0049$ & 0.8723 \\
2.0 & $0.8790\pm0.0085$ & 0.8734 \\
\bottomrule
\end{tabular}
\end{table}

\subsubsection{组件消融}

表~\ref{tab:component_ablation_zh} 分析两个视图和跨视图对比学习的作用。当前预备记录使用 $k=20$ 的原始有向 KNN 视图以及较早的 $\lambda_{\mathrm{DVCL}}=0.5$ 设置；最终消融实验将与主比较使用相同的冻结划分和攻击产物。

\begin{table}[htbp]
\centering
\caption{ACM 上的组件消融准确率。}
\label{tab:component_ablation_zh}
\begin{tabular}{lcc}
\toprule
变体 & 干净准确率 & 攻击平均准确率 \\
\midrule
完整 DVCL & $0.8770\pm0.0071$ & 0.8710 \\
去除对比损失 & $0.8743\pm0.0074$ & 0.8656 \\
仅拓扑视图 & $0.8523\pm0.0177$ & 0.8471 \\
仅特征视图 & $0.8539\pm0.0086$ & 0.8539 \\
\bottomrule
\end{tabular}
\end{table}

\subsubsection{特征视图构造}

表~\ref{tab:knn_supp_zh} 给出特征视图构造的预备三随机种子比较。该分析属于补充实验，不改变主比较中固定使用的原始有向 KNN 配置。

\begin{table}[htbp]
\centering
\caption{ACM 上三随机种子的特征视图构造准确率。}
\label{tab:knn_supp_zh}
\begin{tabular}{lcc}
\toprule
特征视图 & 干净准确率 & 攻击平均准确率 \\
\midrule
原始 KNN，$k=10$ & 0.8750 & 0.8719 \\
原始 KNN，$k=20$ & 0.8752 & 0.8706 \\
原始 KNN，$k=40$ & 0.8665 & 0.8677 \\
互近邻原始 KNN，$k=20$ & 0.8804 & 0.8739 \\
异构 profile KNN，$k=20$ & 0.8807 & 0.8784 \\
\bottomrule
\end{tabular}
\end{table}

\FloatBarrier
\subsection{附加鲁棒性分析}

表~\ref{tab:adaptive_attack_zh} 预留用于检验：当攻击者知道 DVCL 完整架构时，特征诱导视图是否仍然可靠。在这些结果完成之前，论文不能把结论扩展到结构攻击之外。

\begin{table}[htbp]
\centering
\small
\caption{超出纯拓扑扰动范围的攻击准确率。}
\label{tab:adaptive_attack_zh}
\begin{tabular}{lcccc}
\toprule
攻击设置 & HSeCo$^{\dagger}$ & RoHe & HeteGuard & DVCL \\
\midrule
特征扰动 & -- & -- & -- & -- \\
拓扑与特征联合攻击 & -- & -- & -- & -- \\
双视图自适应攻击 & -- & -- & -- & -- \\
\bottomrule
\end{tabular}
\end{table}

\FloatBarrier
\subsection{可复现性与结论边界}

在最终比较前，本文将修复 HSeCo$^{\dagger}$ 的早停逻辑，同时保存和恢复语义模块与分类器，并基于恢复后的 checkpoint 重新生成语义图。对于每个数据集和随机种子，本文保存包含划分索引、攻击产物路径、目标节点 ID、随机种子、模型配置和所选 checkpoint 的 manifest。之后，所有基线都从这些冻结 manifest 重新运行。

当前阶段只能支持一个受限结论：在统一 ACM 协议下，DVCL 的预备结果在 PRBCD、HetePRBCD 和 FGSM 风格结构目标攻击下优于当前 HSeCo 复现。要进一步声称 DVCL 优于现有鲁棒异构图方法，必须补齐多个数据集以及表~\ref{tab:baselines_zh} 中外部基线的结果。

\FloatBarrier
